\documentclass[10pt,a4paper]{amsart}
\usepackage{amsmath,amssymb,amsthm,graphicx}
\usepackage{verbatim}
\newenvironment{reference}{\comment}{\endcomment}
%\newenvironment{reference}{}{}
\newenvironment{slogan}{\comment}{\endcomment}
\newenvironment{history}{\comment}{\endcomment}

% For commutative diagrams we use Xy-pic
\usepackage[all]{xy}

% We use 2cell for 2-commutative diagrams.
\xyoption{2cell}
% Theorem environments.
%
\theoremstyle{plain}
\newtheorem{theorem}[subsection]{Theorem}
\newtheorem{proposition}[subsection]{Proposition}
\newtheorem{lemma}[subsection]{Lemma}

\theoremstyle{definition}
\newtheorem{definition}[subsection]{Definition}
\newtheorem{example}[subsection]{Example}
\newtheorem{exercise}[subsection]{Exercise}
\newtheorem{situation}[subsection]{Situation}

\theoremstyle{remark}
\newtheorem{remark}[subsection]{Remark}
\newtheorem{remarks}[subsection]{Remarks}

\numberwithin{equation}{subsection}

% Macros
%
\def\lim{\mathop{\mathrm{lim}}\nolimits}
\def\colim{\mathop{\mathrm{colim}}\nolimits}
\def\Spec{\mathop{\mathrm{Spec}}}
\def\Hom{\mathop{\mathrm{Hom}}\nolimits}
\def\Ext{\mathop{\mathrm{Ext}}\nolimits}
\def\SheafHom{\mathop{\mathcal{H}\!\mathit{om}}\nolimits}
\def\SheafExt{\mathop{\mathcal{E}\!\mathit{xt}}\nolimits}
\def\Sch{\mathit{Sch}}
\def\Mor{\mathop{\mathrm{Mor}}\nolimits}
\def\Ob{\mathop{\mathrm{Ob}}\nolimits}
\def\Sh{\mathop{\mathit{Sh}}\nolimits}
\def\NL{\mathop{N\!L}\nolimits}
\def\CH{\mathop{\mathrm{CH}}\nolimits}
\def\proetale{{pro\text{-}\acute{e}tale}}
\def\etale{{\acute{e}tale}}
\def\QCoh{\mathit{QCoh}}
\def\Ker{\mathop{\mathrm{Ker}}}
\def\Im{\mathop{\mathrm{Im}}}
\def\Coker{\mathop{\mathrm{Coker}}}
\def\Coim{\mathop{\mathrm{Coim}}}

% Boxtimes
%
\DeclareMathSymbol{\boxtimes}{\mathbin}{AMSa}{"02}

%
% Macros for moduli stacks/spaces
%
\def\QCohstack{\mathcal{QC}\!\mathit{oh}}
\def\Cohstack{\mathcal{C}\!\mathit{oh}}
\def\Spacesstack{\mathcal{S}\!\mathit{paces}}
\def\Quotfunctor{\mathrm{Quot}}
\def\Hilbfunctor{\mathrm{Hilb}}
\def\Curvesstack{\mathcal{C}\!\mathit{urves}}
\def\Polarizedstack{\mathcal{P}\!\mathit{olarized}}
\def\Complexesstack{\mathcal{C}\!\mathit{omplexes}}
% \Pic is the operator that assigns to X its picard group, usage \Pic(X)
% \Picardstack_{X/B} denotes the Picard stack of X over B
% \Picardfunctor_{X/B} denotes the Picard functor of X over B
\def\Pic{\mathop{\mathrm{Pic}}\nolimits}
\def\Picardstack{\mathcal{P}\!\mathit{ic}}
\def\Picardfunctor{\mathrm{Pic}}
\def\Deformationcategory{\mathcal{D}\!\mathit{ef}}
\title{Dichotomy for irreducible polynomials: coprime derivative or p-power substitution}
\author{The Stacks Project Authors}
\date{stacks-project:fields.tex:local-snapshot-20261012; GNU FDL 1.2 or later}
\begin{document}
\maketitle
\noindent\textit{Editorial note.} All delivered mathematical text is the author's own, copied from the pinned Stacks Project chapter file fields.tex; no statement or proof step was rewritten, reordered or removed.\par\medskip
\noindent\textit{Editorial note.} The counted proof is the author's complete argument: This lemma is the engine of the purely inseparable theory: P(x) = Q(x\textasciicircum{}q) is what makes alpha\textasciicircum{}q separable in lemma-separable-first.\par\medskip
\noindent\textit{Editorial note.} The label lemma-irreducible-polynomials names the counted statement, and the proof environment immediately following it in the source is delivered verbatim as its proof.\par\medskip
\noindent\textit{Editorial note.} The counted statement is carried by the author's own lemma environment in the chapter source.\par\medskip
\noindent\textit{Allowed standard prerequisites (editorial).}\begin{itemize}
\item Polynomial algebra over a field: k[x] is a PID so (P, P') = (R); the derivative of a polynomial in characteristic p kills exactly the terms of degree divisible by p
\end{itemize}\medskip
\begin{lemma}
\label{lemma-irreducible-polynomials}
Let $F$ be a field. Let $P \in F[x]$ be an irreducible polynomial over $F$.
Let $P' = \text{d}P/\text{d}x$ be the derivative of $P$ with respect
to $x$. Then one of the following two cases happens
\begin{enumerate}
\item $P$ and $P'$ are relatively prime, or
\item $P'$ is the zero polynomial.
\end{enumerate}
The second case can only happen if $F$ has characteristic $p > 0$.
In this case $P(x) = Q(x^q)$ where $q = p^f$ is a power of $p$ and
$Q \in F[x]$ is an irreducible polynomial such that $Q$ and $Q'$
are relatively prime.
\end{lemma}

\begin{proof}
Note that $P'$ has degree $< \deg(P)$. Hence if $P$ and $P'$ are not relatively
prime, then $(P, P') = (R)$ where $R$ is a polynomial of degree $< \deg(P)$
contradicting the irreducibility of $P$. This proves we have the dichotomy
between (1) and (2).

\medskip\noindent
Assume we are in case (2) and $P = a_d x^d + \ldots + a_0$. Then
$P' = da_d x^{d - 1} + \ldots + a_1$. In characteristic $0$ we see
that this forces $a_d, \ldots, a_1 = 0$ which would mean $P$ is constant
a contradiction. Thus we conclude that the characteristic $p$ is positive.
In this case the condition $P' = 0$ forces $a_i = 0$ whenever $p$ does
not divide $i$.
In other words, $P(x) = P_1(x^p)$ for some nonconstant polynomial $P_1$.
Clearly, $P_1$ is irreducible as well. By induction on the degree we
see that $P_1(x) = Q(x^q)$ as in the statement of the lemma, hence
$P(x) = Q(x^{pq})$ and the lemma is proved.
\end{proof}

\end{document}
